\documentclass[11pt]{article}
\usepackage[margin=25mm]{geometry}
\usepackage{fontspec}
\setmainfont{Latin Modern Roman}
\usepackage{amsmath,amssymb,mathtools}
\usepackage{unicode-math}
\setmathfont{Latin Modern Math}
\usepackage{xurl,hyperref}
\hypersetup{colorlinks=true,urlcolor=blue}
\setcounter{secnumdepth}{0}
\setlength{\parindent}{0pt}
\setlength{\parskip}{5pt}
\setlength{\emergencystretch}{3em}
\newcommand{\sref}[2]{\hyperlink{source:#1:#2}{[S#2]}}
\newcommand{\cataloguescope}[1]{\paragraph{Recorded target.}#1}
\newcommand{\R}{\mathbb R}
\newcommand{\C}{\mathbb C}
\newcommand{\N}{\mathbb N}
\newcommand{\Z}{\mathbb Z}
\newcommand{\Q}{\mathbb Q}
\newcommand{\F}{\mathbb F}
\newcommand{\eps}{\varepsilon}
\begin{document}
\section{Full-space mass-splitting criterion for prograde relative periodic orbits}\label{statement}

\subsection{Definitions and mathematical statement}

\paragraph{Shared notation.}$\mathbb N=\{1,2,\ldots\}$; the standard number fields have their usual meanings. All additional conventions are stated locally.


\paragraph{Definitions and notation.}In the Newtonian $n$-body problem, positions $q_i\in\mathbb R^d$ and masses $m_i>0$ obey $m_i\ddot q_i=\nabla_{q_i}U$, where $U(q)=\sum_{i<j}m_im_j/|q_i-q_j|$ and the gravitational constant is normalized to one. Here $T=\tfrac12\sum_i m_i|\dot q_i|^2$, $E=T-U$, and a collision has $q_i=q_j$. The center of mass can be fixed at zero; moment of inertia is $I=\sum_i m_i|q_i|^2$. A prograde three-body motion has the relevant orbital rotations in the same orientation under a chosen angular-momentum or Jacobi-coordinate convention. Comparable masses means mass ratios stay away from a zero-mass perturbative limit. Avoiding a nearly zero mutual-distance ratio excludes a hierarchical binary-plus-distant-body perturbative regime.

\paragraph{Catalogue wording retained for traceability.}
Given three distinct masses comparable in size, prove existence of prograde orbits without assuming that one ratio of mutual distances is nearly zero.
\paragraph{Scope and source qualification.}The source records continuation arguments in small-mass or hierarchical regimes and an equal-two-mass case. It does not provide numerical comparability bounds, a prograde convention or a specified orbit family; these require explicit formulation in a proof. The catalogue title/status was paired with a machine-generated research attempt; that attempt is not peer-reviewed evidence of settlement. The statement covers the original source questions rather than only the title’s subcase.
\subsection{Short English statement}
Prove nonhierarchical prograde three-body motions for three distinct comparable masses.
\subsection{Research attempt}
\paragraph{Model, date and investigated scope.}
OpenAI GPT-6 Astra (configured model \texttt{gpt-6-astra}), ultra reasoning; generated 11 October 2026. Proved conditional preservation of collision clearance, nonhierarchy, prograde margins and comparability along a nondegenerate shooting branch. No seed certificate or all-comparable-mass existence theorem.

\paragraph{Sources and literature check.}
Read the supplied complete definition, catalogue wording, scope qualifications and source list for this canonical ID. The original-source citations are inherited as supplied references; no fresh primary-source or novelty search was performed for this attempt. No current open-status claim is inferred from this calculation.

\paragraph{A quantitative openness statement for a hypothetical seed.}Fix a collision-free relative-periodic solution at a positive mass triple \(m^0\), and assume a smooth finite-dimensional shooting map \(F(x,m)=0\) has been formed after fixing time/rotation freedoms. Suppose \(F(x_0,m^0)=0\) and \(D_xF\) is invertible. The implicit-function theorem then gives a unique nearby branch \(x(m)\). Smooth dependence of collision-free ODE solutions on parameters yields uniform convergence of the corresponding orbit segments on their common compact time interval.
If the seed has minimum mutual distance \(d_*>0\), maximum distance \(D_*<\infty\), and a strict positive prograde diagnostic \(P_*>0\), continuity preserves, for sufficiently small mass changes,
\[d_{\min}\ge d_*/2,\quad d_{\max}\le2D_*,\quad d_{\min}/d_{\max}\ge d_*/(4D_*),\quad P\ge P_*/2.
\]
Here P must be a specified continuous angular-momentum/Jacobi-rotation diagnostic on the full orbit; pointwise nonnegative values without a strict margin would not suffice. Any strict mass-comparability bounds likewise persist. A neighborhood of an equal-pair seed contains all-distinct positive triples.
\paragraph{The missing certificate.}The argument assumes exactly what must be verified in the intended prograde family: a seed with strict nonhierarchical and prograde margins and an invertible fully reduced shooting derivative. Invertibility in a smaller symmetry-restricted space would not justify continuation under mass perturbations breaking that symmetry. No such seed or derivative is computed here.
\paragraph{Scope.}Even a verified local branch would cover only nearby distinct masses, not every arbitrary comparable triple. The source's intended orbit family and prograde convention must also be fixed. The first gap is the concrete nondegeneracy certificate; extending across a large mass domain would require additional global continuation and exclusion of collisions or bifurcation loss. This is a standard conditional continuation argument, not an unconditional existence theorem.

\paragraph{Verification and adversarial audit.}
Checked strict margins, compact-time parameter continuity, distance-ratio constant, all-distinct nearby triples and full versus symmetry-restricted nondegeneracy. This was an LLM adversarial audit, not independent human verification; no theorem prover was used. Prompts A then B applied to this ID: Checked strict margins, compact-time parameter continuity, distance-ratio constant, all-distinct nearby triples and full versus symmetry-restricted nondegeneracy. Scope remains limited: Verify a full-space reduced shooting nondegeneracy certificate for an actual nonhierarchical prograde seed, then control global continuation if all masses are required.

Coordinator checked the conditional continuation margins and the need for nondegeneracy after full symmetry reduction; no seed existence is inferred. This additional check was model-assisted, not independent human verification.

\paragraph{Tools, process and contributions.}
Python 3 and local shell used to read supplied files and serialize this manuscript; mathematics derived symbolically in the text, with no numerical experiment or formal proof assistant. Contribution prompt A was followed by prompt B. GPT-6 Astra with ultra reasoning drafted the research and performed its stated self-audit; any additional coordinator contribution is identified in the verification and audit. A GPT-6-based Codex coordinator assembled the manuscript and checked repository integration; its exact variant and reasoning configuration were not exposed. No human mathematical contribution, independent human proof check or formal proof certification is supplied. Complete derivations and the consulted source references are included here; private reasoning transcripts are not published.

\paragraph{Final status and first remaining gap.}
\textbf{UNRESOLVED}. Verify a full-space reduced shooting nondegeneracy certificate for an actual nonhierarchical prograde seed, then control global continuation if all masses are required.

\section{Completion Estimate}
\noindent\textbf{COMPLETION: 0\%}
The calculation is an elementary benchmark or conditional reduction; it establishes no claimed new portion of the complete source target.

\subsection{Sources}
\begin{itemize}
\item[\hypertarget{source:statement:1}{S1}] ULAM.AI and the UnsolvedMath Contributors, UnsolvedMath, catalogue entry. Catalogue curation under CC BY 4.0; definitions and English exposition added locally.
\url{https://huggingface.co/datasets/ulamai/UnsolvedMath}.
\item[\hypertarget{source:statement:2}{S2}] American Institute of Mathematics, Variational Methods in Celestial Mechanics; notes by Kuo-Chang Chen, version 22 June 2003, source item 29.
\url{https://aimath.org/WWN/varcelest/varcelest.pdf}.
\end{itemize}
\label{end:statement}
\end{document}
